\section{Panduan membaca tabel lengkap}\label{app:read}
Lampiran mempertahankan 176 model utama dan 340 sensitivitas, termasuk model yang belum mendukung level atau mempunyai catatan diagnostik. Kolom koefisien menggunakan $\beta_q$ pada ARDL dan $\beta_-$ pada NARDL. Dengan definisi kurs yang digunakan, koefisien positif mengarah pada perbaikan rasio neraca ketika depresiasi, sedangkan negatif mengarah pada pemburukan. Penafsiran itu tetap memerlukan kelayakan hubungan level.

\subsection{Kode pasar, kelompok, dan varian}
Negara pelapor semua tabel adalah Indonesia. \texttt{WLD} berarti seluruh negara mitra; \texttt{AUS}, \texttt{CHN}, \texttt{IND}, \texttt{JPN}, \texttt{KOR}, \texttt{MYS}, \texttt{SGP}, \texttt{THA}, \texttt{USA}, dan \texttt{VNM} adalah mitra yang sudah dibahas. G0 berarti seluruh barang pada cakupan pasar tersebut. G1--G7 mengikuti Tabel~\ref{tab:groups}.

Varian U adalah baseline maksimum empat lag level; L2 maksimum dua; P tambahan kontrol BCOM; W06 perluasan window agregat Indonesia 2006+; WPI memakai proksi harga WPI India; FX memakai keluarga kurs GEM alternatif. Perbedaan varian dapat melibatkan perubahan window atau konsep sumber. Tabel tidak menjadikan varian-varian tersebut sebagai eksperimen yang hanya berbeda satu unsur.

\subsection{Arti status}
\texttt{L+D} berarti syarat level dan penyaringan diagnostik yang digunakan terpenuhi. \texttt{L} berarti syarat level terpenuhi tetapi ada catatan diagnostik. Tanda -- berarti dukungan level lengkap belum terpenuhi atau inferensinya tidak tersedia. Model dengan tanda -- tetap dapat mempunyai koefisien nominal signifikan. Status tidak secara otomatis menentukan bahwa suatu efek berbeda dari nol.

Probabilitas BH hanya dihitung pada keluarga model utama. Kolom BH sensitivitas karena itu ditulis --. Probabilitas lebih kecil dari 0,001 dicetak sebagai $<0{,}001$, bukan sebagai nol persis. Jumlah observasi adalah sampel estimasi setelah lag, bukan jumlah bulan atau kuartal raw.

\subsection{Berkas yang memuat rincian lebih jauh}
Workbook \path{output/tables/oskooee/Oskooee_Results.xlsx} memuat seluruh koefisien, unit root, registry sampel dan status. \path{results_all_models.csv} serta \path{diagnostic_summary.csv} menyimpan probabilitas masing-masing diagnostik. Identitas model lengkap pada \path{assets/all_516_models_source.csv} memungkinkan tabel ringkas ini dipadankan dengan hasil Stata.

Atlas respons lengkap berada dalam PPT versi lengkap dan dipetakan melalui \path{model_response_slide_index.csv} pada folder presentasi. Paper memuat grafik terpilih untuk pendalaman, sedangkan tabel berikut mempertahankan semua estimasi numeriknya.

\clearpage\section{Seluruh 176 model utama}
\noindent Negara pelapor: Indonesia. Mitra \texttt{WLD} adalah seluruh negara. Kelompok dan status mengikuti Lampiran~\ref{app:read}. Koefisien belum menjadi hubungan jangka panjang yang mapan pada baris tanpa dukungan level.
{\footnotesize\setlength{\tabcolsep}{3pt}\renewcommand{\arraystretch}{1.15}
\begin{longtable}{p{0.075\linewidth}p{0.05\linewidth}p{0.09\linewidth}p{0.04\linewidth}p{0.05\linewidth}p{0.1\linewidth}p{0.085\linewidth}p{0.085\linewidth}p{0.075\linewidth}}
\toprule
\textbf{Mitra} & \textbf{Kel.} & \textbf{Model} & \textbf{N} & \textbf{Var.} & \textbf{Koef.} & \textbf{p} & \textbf{BH} & \textbf{Status}\\
\midrule
\endfirsthead
\toprule
\textbf{Mitra} & \textbf{Kel.} & \textbf{Model} & \textbf{N} & \textbf{Var.} & \textbf{Koef.} & \textbf{p} & \textbf{BH} & \textbf{Status}\\
\midrule
\endhead
\midrule\multicolumn{9}{r}{\footnotesize Bersambung ke halaman berikutnya}\\\endfoot
\bottomrule\endlastfoot
AUS & G3 & ARDL & 77 & U & -1,770 & 0,043 & 0,165 & --\\
AUS & G3 & NARDL & 77 & U & -3,345 & \(<0{,}001\) & 0,004 & L\\
AUS & G5 & ARDL & 77 & U & 4,002 & 0,120 & 0,307 & --\\
AUS & G5 & NARDL & 77 & U & 7,466 & 0,041 & 0,138 & --\\
AUS & G6 & ARDL & 77 & U & 1,701 & 0,008 & 0,103 & L+D\\
AUS & G6 & NARDL & 77 & U & 1,379 & 0,142 & 0,332 & L+D\\
AUS & G2 & ARDL & 77 & U & 0,704 & 0,316 & 0,544 & --\\
AUS & G2 & NARDL & 77 & U & 1,088 & 0,285 & 0,488 & --\\
AUS & G7 & ARDL & 77 & U & 3,265 & 0,026 & 0,121 & L\\
AUS & G7 & NARDL & 77 & U & 5,725 & 0,006 & 0,036 & L\\
AUS & G4 & ARDL & 77 & U & 1,608 & 0,024 & 0,121 & L\\
AUS & G4 & NARDL & 77 & U & 2,518 & 0,010 & 0,053 & L+D\\
AUS & G1 & ARDL & 77 & U & 0,919 & 0,460 & 0,656 & --\\
AUS & G1 & NARDL & 77 & U & 5,628 & 0,003 & 0,023 & --\\
AUS & G0 & ARDL & 77 & U & 2,084 & 0,002 & 0,021 & L+D\\
AUS & G0 & NARDL & 77 & U & 3,356 & \(<0{,}001\) & 0,004 & L+D\\
CHN & G3 & ARDL & 78 & U & -0,809 & 0,119 & 0,307 & L+D\\
CHN & G3 & NARDL & 78 & U & -2,036 & 0,002 & 0,015 & L+D\\
CHN & G5 & ARDL & 78 & U & 1,211 & 0,325 & 0,544 & L\\
CHN & G5 & NARDL & 78 & U & 0,595 & 0,764 & 0,842 & L\\
CHN & G6 & ARDL & 78 & U & -2,216 & 0,005 & 0,103 & L\\
CHN & G6 & NARDL & 78 & U & -3,458 & \(<0{,}001\) & 0,003 & L\\
CHN & G2 & ARDL & 78 & U & -0,569 & 0,567 & 0,752 & --\\
CHN & G2 & NARDL & 78 & U & -2,711 & 0,031 & 0,108 & L+D\\
CHN & G7 & ARDL & 78 & U & -1,924 & 0,101 & 0,300 & --\\
CHN & G7 & NARDL & 78 & U & -3,658 & 0,022 & 0,101 & L\\
CHN & G4 & ARDL & 78 & U & -3,291 & 0,046 & 0,168 & --\\
CHN & G4 & NARDL & 78 & U & -6,234 & 0,004 & 0,023 & --\\
CHN & G1 & ARDL & 78 & U & 0,300 & 0,878 & 0,955 & --\\
CHN & G1 & NARDL & 78 & U & 3,033 & 0,663 & 0,763 & --\\
CHN & G0 & ARDL & 78 & U & -1,886 & 0,059 & 0,148 & --\\
CHN & G0 & NARDL & 78 & U & -2,727 & 0,250 & 0,350 & --\\
IND & G3 & ARDL & 49 & U & 6,325 & 0,645 & 0,823 & --\\
IND & G3 & NARDL & 49 & U & -5,076 & 0,316 & 0,503 & L\\
IND & G5 & ARDL & 49 & U & 1,012 & 0,913 & 0,963 & --\\
IND & G5 & NARDL & 49 & U & 0,815 & 0,939 & 0,970 & --\\
IND & G6 & ARDL & 49 & U & 5,476 & 0,549 & 0,742 & --\\
IND & G6 & NARDL & 49 & U & -4,569 & 0,320 & 0,503 & L+D\\
IND & G2 & ARDL & 49 & U & 3,982 & 0,535 & 0,736 & --\\
IND & G2 & NARDL & 49 & U & -8,904 & 0,090 & 0,263 & L\\
IND & G7 & ARDL & 49 & U & 7,828 & 0,358 & 0,574 & --\\
IND & G7 & NARDL & 49 & U & -4,506 & 0,346 & 0,526 & L+D\\
IND & G4 & ARDL & 49 & U & 10,846 & 0,240 & 0,463 & --\\
IND & G4 & NARDL & 49 & U & 2,938 & 0,702 & 0,795 & --\\
IND & G1 & ARDL & 49 & U & -4,502 & 0,393 & 0,606 & --\\
IND & G1 & NARDL & 49 & U & -6,379 & 0,179 & 0,377 & --\\
IND & G0 & ARDL & 49 & U & 5,254 & 0,437 & 0,486 & --\\
IND & G0 & NARDL & 49 & U & -6,443 & 0,236 & 0,350 & --\\
JPN & G3 & ARDL & 78 & U & -0,075 & 0,925 & 0,963 & --\\
JPN & G3 & NARDL & 78 & U & -1,712 & 0,009 & 0,047 & --\\
JPN & G5 & ARDL & 78 & U & -3,068 & 0,075 & 0,241 & --\\
JPN & G5 & NARDL & 78 & U & -4,687 & 0,002 & 0,018 & --\\
JPN & G6 & ARDL & 78 & U & -1,771 & 0,009 & 0,103 & L\\
JPN & G6 & NARDL & 78 & U & -0,710 & 0,493 & 0,612 & L+D\\
JPN & G2 & ARDL & 78 & U & -0,718 & 0,337 & 0,553 & L\\
JPN & G2 & NARDL & 78 & U & -1,246 & 0,244 & 0,448 & L\\
JPN & G7 & ARDL & 78 & U & -2,107 & 0,016 & 0,119 & L\\
JPN & G7 & NARDL & 78 & U & -0,160 & 0,894 & 0,945 & L+D\\
JPN & G4 & ARDL & 78 & U & -2,347 & 0,021 & 0,121 & --\\
JPN & G4 & NARDL & 78 & U & -3,048 & 0,002 & 0,015 & --\\
JPN & G1 & ARDL & 78 & U & -1,171 & 0,166 & 0,375 & --\\
JPN & G1 & NARDL & 78 & U & -1,477 & 0,159 & 0,350 & --\\
JPN & G0 & ARDL & 78 & U & -2,106 & 0,017 & 0,058 & --\\
JPN & G0 & NARDL & 78 & U & -2,456 & 0,006 & 0,020 & --\\
KOR & G3 & ARDL & 78 & U & 1,449 & 0,056 & 0,187 & --\\
KOR & G3 & NARDL & 78 & U & 1,601 & 0,241 & 0,448 & --\\
KOR & G5 & ARDL & 78 & U & 1,416 & 0,420 & 0,622 & --\\
KOR & G5 & NARDL & 78 & U & 7,147 & 0,027 & 0,104 & L+D\\
KOR & G6 & ARDL & 78 & U & -0,370 & 0,857 & 0,955 & --\\
KOR & G6 & NARDL & 78 & U & 4,662 & 0,092 & 0,263 & --\\
KOR & G2 & ARDL & 78 & U & 7,941 & 0,118 & 0,307 & --\\
KOR & G2 & NARDL & 78 & U & 0,875 & 0,945 & 0,970 & --\\
KOR & G7 & ARDL & 78 & U & -0,712 & 0,773 & 0,910 & --\\
KOR & G7 & NARDL & 78 & U & 5,815 & 0,109 & 0,279 & --\\
KOR & G4 & ARDL & 78 & U & 0,887 & 0,652 & 0,823 & --\\
KOR & G4 & NARDL & 78 & U & 10,950 & 0,102 & 0,275 & --\\
KOR & G1 & ARDL & 78 & U & 1,747 & 0,055 & 0,187 & --\\
KOR & G1 & NARDL & 78 & U & 4,717 & 0,001 & 0,013 & --\\
KOR & G0 & ARDL & 78 & U & -0,074 & 0,944 & 0,944 & --\\
KOR & G0 & NARDL & 78 & U & 3,935 & 0,144 & 0,288 & --\\
MYS & G3 & ARDL & 78 & U & 0,338 & 0,803 & 0,922 & --\\
MYS & G3 & NARDL & 78 & U & 3,464 & 0,181 & 0,377 & --\\
MYS & G5 & ARDL & 78 & U & 0,178 & 0,751 & 0,904 & L+D\\
MYS & G5 & NARDL & 78 & U & 0,784 & 0,468 & 0,610 & L+D\\
MYS & G6 & ARDL & 78 & U & -1,290 & 0,088 & 0,270 & L+D\\
MYS & G6 & NARDL & 78 & U & 1,444 & 0,258 & 0,461 & L+D\\
MYS & G2 & ARDL & 78 & U & -0,746 & 0,695 & 0,850 & --\\
MYS & G2 & NARDL & 78 & U & 4,491 & 0,068 & 0,219 & --\\
MYS & G7 & ARDL & 78 & U & -1,376 & 0,116 & 0,307 & L+D\\
MYS & G7 & NARDL & 78 & U & 1,117 & 0,483 & 0,610 & L+D\\
MYS & G4 & ARDL & 78 & U & -0,645 & 0,408 & 0,617 & --\\
MYS & G4 & NARDL & 78 & U & 1,176 & 0,309 & 0,503 & --\\
MYS & G1 & ARDL & 78 & U & 0,285 & 0,845 & 0,955 & --\\
MYS & G1 & NARDL & 78 & U & 1,939 & 0,406 & 0,578 & --\\
MYS & G0 & ARDL & 78 & U & -0,716 & 0,081 & 0,163 & L\\
MYS & G0 & NARDL & 78 & U & 1,346 & 0,024 & 0,059 & L+D\\
SGP & G3 & ARDL & 78 & U & 5,982 & 0,027 & 0,121 & --\\
SGP & G3 & NARDL & 78 & U & 4,941 & 0,092 & 0,263 & --\\
SGP & G5 & ARDL & 78 & U & -0,469 & 0,669 & 0,831 & --\\
SGP & G5 & NARDL & 78 & U & -1,303 & 0,470 & 0,610 & --\\
SGP & G6 & ARDL & 78 & U & -1,192 & 0,025 & 0,121 & --\\
SGP & G6 & NARDL & 78 & U & -0,701 & 0,483 & 0,610 & --\\
SGP & G2 & ARDL & 78 & U & 2,616 & 0,325 & 0,544 & L\\
SGP & G2 & NARDL & 78 & U & -0,102 & 0,985 & 0,991 & --\\
SGP & G7 & ARDL & 78 & U & -1,383 & 0,016 & 0,119 & --\\
SGP & G7 & NARDL & 78 & U & -0,709 & 0,504 & 0,617 & --\\
SGP & G4 & ARDL & 78 & U & 0,424 & 0,604 & 0,788 & --\\
SGP & G4 & NARDL & 78 & U & -2,032 & 0,022 & 0,101 & L\\
SGP & G1 & ARDL & 78 & U & 0,007 & 0,994 & 0,994 & --\\
SGP & G1 & NARDL & 78 & U & -1,544 & 0,308 & 0,503 & --\\
SGP & G0 & ARDL & 78 & U & -0,451 & 0,432 & 0,486 & --\\
SGP & G0 & NARDL & 78 & U & -0,262 & 0,810 & 0,810 & --\\
THA & G3 & ARDL & 78 & U & 1,906 & 0,216 & 0,442 & --\\
THA & G3 & NARDL & 78 & U & 3,848 & 0,155 & 0,350 & L\\
THA & G5 & ARDL & 78 & U & -1,027 & 0,433 & 0,629 & --\\
THA & G5 & NARDL & 78 & U & 2,593 & 0,355 & 0,526 & --\\
THA & G6 & ARDL & 78 & U & 1,047 & 0,137 & 0,341 & --\\
THA & G6 & NARDL & 78 & U & 3,454 & 0,024 & 0,104 & --\\
THA & G2 & ARDL & 78 & U & 2,543 & 0,015 & 0,119 & L\\
THA & G2 & NARDL & 78 & U & 0,354 & 0,896 & 0,945 & --\\
THA & G7 & ARDL & 78 & U & -0,592 & 0,270 & 0,496 & --\\
THA & G7 & NARDL & 78 & U & 0,755 & 0,594 & 0,703 & --\\
THA & G4 & ARDL & 78 & U & 1,063 & 0,031 & 0,124 & L+D\\
THA & G4 & NARDL & 78 & U & -1,568 & 0,124 & 0,309 & L+D\\
THA & G1 & ARDL & 78 & U & 1,001 & 0,218 & 0,442 & --\\
THA & G1 & NARDL & 78 & U & 0,564 & 0,766 & 0,842 & --\\
THA & G0 & ARDL & 78 & U & 1,180 & 0,006 & 0,031 & L\\
THA & G0 & NARDL & 78 & U & 0,950 & 0,280 & 0,350 & L\\
USA & G3 & ARDL & 75 & U & -0,836 & 0,017 & 0,119 & L+D\\
USA & G3 & NARDL & 75 & U & 0,622 & 0,278 & 0,487 & L+D\\
USA & G5 & ARDL & 75 & U & 1,166 & 0,780 & 0,910 & --\\
USA & G5 & NARDL & 75 & U & 4,530 & 0,434 & 0,587 & --\\
USA & G6 & ARDL & 75 & U & -0,088 & 0,912 & 0,963 & --\\
USA & G6 & NARDL & 75 & U & 3,592 & \(<0{,}001\) & 0,008 & L\\
USA & G2 & ARDL & 75 & U & -0,335 & 0,880 & 0,955 & --\\
USA & G2 & NARDL & 75 & U & 5,861 & \(<0{,}001\) & 0,003 & L+D\\
USA & G7 & ARDL & 75 & U & 5,869 & 0,214 & 0,442 & --\\
USA & G7 & NARDL & 75 & U & 13,527 & 0,354 & 0,526 & --\\
USA & G4 & ARDL & 75 & U & -0,313 & 0,517 & 0,723 & L\\
USA & G4 & NARDL & 75 & U & 2,924 & \(<0{,}001\) & 0,010 & L\\
USA & G1 & ARDL & 75 & U & 1,517 & 0,196 & 0,430 & --\\
USA & G1 & NARDL & 75 & U & 2,189 & 0,142 & 0,332 & --\\
USA & G0 & ARDL & 75 & U & 0,989 & 0,163 & 0,272 & --\\
USA & G0 & NARDL & 75 & U & 3,095 & 0,006 & 0,020 & --\\
VNM & G3 & ARDL & 74 & U & 5,911 & 0,009 & 0,103 & --\\
VNM & G3 & NARDL & 74 & U & 7,260 & 0,026 & 0,104 & --\\
VNM & G5 & ARDL & 74 & U & 2,773 & 0,148 & 0,356 & --\\
VNM & G5 & NARDL & 74 & U & -0,961 & 0,621 & 0,725 & L\\
VNM & G6 & ARDL & 74 & U & 2,123 & 0,252 & 0,474 & --\\
VNM & G6 & NARDL & 74 & U & 2,287 & 0,429 & 0,587 & --\\
VNM & G2 & ARDL & 74 & U & 3,995 & 0,006 & 0,103 & L\\
VNM & G2 & NARDL & 74 & U & 4,684 & 0,031 & 0,108 & L\\
VNM & G7 & ARDL & 74 & U & -3,966 & 0,019 & 0,121 & L\\
VNM & G7 & NARDL & 74 & U & -4,206 & 0,103 & 0,275 & L\\
VNM & G4 & ARDL & 74 & U & -1,282 & 0,158 & 0,368 & L\\
VNM & G4 & NARDL & 74 & U & -2,069 & 0,186 & 0,377 & L\\
VNM & G1 & ARDL & 74 & U & 2,785 & 0,300 & 0,536 & --\\
VNM & G1 & NARDL & 74 & U & -0,025 & 0,991 & 0,991 & --\\
VNM & G0 & ARDL & 74 & U & 1,389 & 0,205 & 0,293 & --\\
VNM & G0 & NARDL & 74 & U & -0,551 & 0,737 & 0,810 & --\\
WLD & G3 & ARDL & 61 & U & 0,005 & 0,993 & 0,994 & --\\
WLD & G3 & NARDL & 61 & U & -1,621 & 0,192 & 0,380 & --\\
WLD & G5 & ARDL & 61 & U & 1,425 & 0,373 & 0,586 & --\\
WLD & G5 & NARDL & 61 & U & 1,589 & 0,402 & 0,578 & --\\
WLD & G6 & ARDL & 61 & U & -1,355 & \(<0{,}001\) & 0,004 & L+D\\
WLD & G6 & NARDL & 61 & U & -0,504 & 0,539 & 0,649 & --\\
WLD & G2 & ARDL & 61 & U & 2,026 & 0,031 & 0,124 & L+D\\
WLD & G2 & NARDL & 61 & U & 4,580 & \(<0{,}001\) & 0,004 & L\\
WLD & G7 & ARDL & 61 & U & -2,040 & \(<0{,}001\) & \(<0{,}001\) & L+D\\
WLD & G7 & NARDL & 61 & U & -0,112 & 0,827 & 0,897 & L\\
WLD & G4 & ARDL & 61 & U & -0,139 & 0,967 & 0,993 & --\\
WLD & G4 & NARDL & 61 & U & 7,851 & 0,431 & 0,587 & --\\
WLD & G1 & ARDL & 61 & U & 0,865 & 0,241 & 0,463 & --\\
WLD & G1 & NARDL & 61 & U & 1,358 & 0,218 & 0,420 & --\\
WLD & G0 & ARDL & 61 & U & 0,429 & 0,661 & 0,661 & --\\
WLD & G0 & NARDL & 61 & U & 2,572 & 0,040 & 0,040 & --\\
\end{longtable}
}

\clearpage\section{Seluruh 340 model sensitivitas}
\noindent Varian dan status mengikuti Lampiran~\ref{app:read}. Estimasi yang berbeda dapat mempunyai window berbeda. BH tidak dihitung untuk keluarga sensitivitas dalam run ini.
{\footnotesize\setlength{\tabcolsep}{3pt}\renewcommand{\arraystretch}{1.15}
\begin{longtable}{p{0.075\linewidth}p{0.05\linewidth}p{0.09\linewidth}p{0.04\linewidth}p{0.05\linewidth}p{0.1\linewidth}p{0.085\linewidth}p{0.085\linewidth}p{0.075\linewidth}}
\toprule
\textbf{Mitra} & \textbf{Kel.} & \textbf{Model} & \textbf{N} & \textbf{Var.} & \textbf{Koef.} & \textbf{p} & \textbf{BH} & \textbf{Status}\\
\midrule
\endfirsthead
\toprule
\textbf{Mitra} & \textbf{Kel.} & \textbf{Model} & \textbf{N} & \textbf{Var.} & \textbf{Koef.} & \textbf{p} & \textbf{BH} & \textbf{Status}\\
\midrule
\endhead
\midrule\multicolumn{9}{r}{\footnotesize Bersambung ke halaman berikutnya}\\\endfoot
\bottomrule\endlastfoot
AUS & G3 & ARDL & 77 & P & -1,803 & 0,153 & -- & --\\
AUS & G3 & NARDL & 77 & P & -3,639 & \(<0{,}001\) & -- & L\\
AUS & G5 & ARDL & 77 & P & 3,972 & 0,266 & -- & --\\
AUS & G5 & NARDL & 77 & P & 6,788 & 0,094 & -- & --\\
AUS & G6 & ARDL & 77 & P & -0,374 & 0,678 & -- & L+D\\
AUS & G6 & NARDL & 77 & P & -0,914 & 0,313 & -- & L\\
AUS & G2 & ARDL & 77 & P & 1,244 & 0,189 & -- & --\\
AUS & G2 & NARDL & 77 & P & 1,141 & 0,292 & -- & --\\
AUS & G7 & ARDL & 77 & P & 2,080 & 0,195 & -- & L+D\\
AUS & G7 & NARDL & 77 & P & 4,137 & 0,022 & -- & L+D\\
AUS & G4 & ARDL & 77 & P & 0,996 & 0,267 & -- & L\\
AUS & G4 & NARDL & 77 & P & 1,618 & 0,122 & -- & L+D\\
AUS & G1 & ARDL & 77 & P & 0,919 & 0,603 & -- & --\\
AUS & G1 & NARDL & 77 & P & 4,610 & 0,029 & -- & --\\
CHN & G3 & ARDL & 78 & P & -0,171 & 0,641 & -- & L+D\\
CHN & G3 & NARDL & 78 & P & -1,336 & 0,020 & -- & L+D\\
CHN & G5 & ARDL & 78 & P & 0,970 & 0,457 & -- & --\\
CHN & G5 & NARDL & 78 & P & -1,667 & 0,418 & -- & L\\
CHN & G6 & ARDL & 78 & P & -1,437 & 0,049 & -- & L\\
CHN & G6 & NARDL & 78 & P & -2,840 & 0,032 & -- & L\\
CHN & G2 & ARDL & 78 & P & 0,211 & 0,844 & -- & --\\
CHN & G2 & NARDL & 78 & P & -2,343 & 0,133 & -- & L+D\\
CHN & G7 & ARDL & 78 & P & -1,692 & 0,204 & -- & --\\
CHN & G7 & NARDL & 78 & P & -3,921 & 0,128 & -- & --\\
CHN & G4 & ARDL & 78 & P & -10,879 & 0,671 & -- & --\\
CHN & G4 & NARDL & 78 & P & -5,276 & 0,069 & -- & --\\
CHN & G1 & ARDL & 78 & P & -0,619 & 0,751 & -- & --\\
CHN & G1 & NARDL & 78 & P & 3,392 & 0,634 & -- & --\\
IND & G3 & ARDL & 49 & P & 5,989 & 0,336 & -- & --\\
IND & G3 & NARDL & 49 & P & 2,367 & 0,708 & -- & --\\
IND & G5 & ARDL & 49 & P & 4,031 & 0,602 & -- & --\\
IND & G5 & NARDL & 49 & P & 4,253 & 0,661 & -- & --\\
IND & G6 & ARDL & 49 & P & 4,372 & 0,086 & -- & L+D\\
IND & G6 & NARDL & 49 & P & -2,073 & 0,506 & -- & --\\
IND & G2 & ARDL & 49 & P & 3,604 & 0,334 & -- & --\\
IND & G2 & NARDL & 49 & P & -4,509 & 0,192 & -- & L\\
IND & G7 & ARDL & 49 & P & 5,995 & 0,303 & -- & --\\
IND & G7 & NARDL & 49 & P & -7,181 & 0,023 & -- & --\\
IND & G4 & ARDL & 49 & P & 15,059 & 0,435 & -- & --\\
IND & G4 & NARDL & 49 & P & 12,149 & 0,614 & -- & --\\
IND & G1 & ARDL & 49 & P & 0,272 & 0,895 & -- & --\\
IND & G1 & NARDL & 49 & P & 0,215 & 0,929 & -- & --\\
JPN & G3 & ARDL & 78 & P & -0,536 & 0,613 & -- & --\\
JPN & G3 & NARDL & 78 & P & -1,996 & 0,026 & -- & --\\
JPN & G5 & ARDL & 78 & P & -3,468 & 0,188 & -- & --\\
JPN & G5 & NARDL & 78 & P & -6,547 & 0,003 & -- & --\\
JPN & G6 & ARDL & 78 & P & -1,534 & 0,025 & -- & L+D\\
JPN & G6 & NARDL & 78 & P & 0,226 & 0,846 & -- & L+D\\
JPN & G2 & ARDL & 78 & P & -0,659 & 0,379 & -- & L\\
JPN & G2 & NARDL & 78 & P & -0,452 & 0,706 & -- & L\\
JPN & G7 & ARDL & 78 & P & -2,007 & 0,022 & -- & L\\
JPN & G7 & NARDL & 78 & P & 0,771 & 0,588 & -- & L+D\\
JPN & G4 & ARDL & 78 & P & -2,522 & 0,037 & -- & --\\
JPN & G4 & NARDL & 78 & P & -2,930 & 0,057 & -- & --\\
JPN & G1 & ARDL & 78 & P & -2,257 & 0,055 & -- & --\\
JPN & G1 & NARDL & 78 & P & -2,745 & 0,034 & -- & --\\
KOR & G3 & ARDL & 78 & P & 1,390 & \(<0{,}001\) & -- & L\\
KOR & G3 & NARDL & 78 & P & 2,348 & \(<0{,}001\) & -- & L\\
KOR & G5 & ARDL & 78 & P & 1,397 & 0,444 & -- & --\\
KOR & G5 & NARDL & 78 & P & 7,249 & 0,019 & -- & --\\
KOR & G6 & ARDL & 78 & P & -0,718 & 0,775 & -- & --\\
KOR & G6 & NARDL & 78 & P & 5,197 & 0,280 & -- & --\\
KOR & G2 & ARDL & 78 & P & 8,328 & 0,078 & -- & --\\
KOR & G2 & NARDL & 78 & P & 2,899 & 0,795 & -- & --\\
KOR & G7 & ARDL & 78 & P & -1,169 & 0,727 & -- & --\\
KOR & G7 & NARDL & 78 & P & 6,791 & 0,302 & -- & --\\
KOR & G4 & ARDL & 78 & P & 0,809 & 0,607 & -- & --\\
KOR & G4 & NARDL & 78 & P & 14,466 & 0,174 & -- & --\\
KOR & G1 & ARDL & 78 & P & 1,704 & 0,074 & -- & --\\
KOR & G1 & NARDL & 78 & P & 4,733 & 0,002 & -- & --\\
MYS & G3 & ARDL & 78 & P & -0,389 & 0,746 & -- & --\\
MYS & G3 & NARDL & 78 & P & 2,669 & 0,120 & -- & --\\
MYS & G5 & ARDL & 78 & P & 0,103 & 0,834 & -- & L+D\\
MYS & G5 & NARDL & 78 & P & 0,341 & 0,736 & -- & L+D\\
MYS & G6 & ARDL & 78 & P & -0,701 & 0,378 & -- & --\\
MYS & G6 & NARDL & 78 & P & 1,152 & 0,371 & -- & L+D\\
MYS & G2 & ARDL & 78 & P & 1,383 & 0,641 & -- & --\\
MYS & G2 & NARDL & 78 & P & 6,257 & 0,011 & -- & --\\
MYS & G7 & ARDL & 78 & P & -0,811 & 0,315 & -- & L+D\\
MYS & G7 & NARDL & 78 & P & 0,462 & 0,760 & -- & L+D\\
MYS & G4 & ARDL & 78 & P & -0,523 & 0,411 & -- & --\\
MYS & G4 & NARDL & 78 & P & 1,546 & 0,069 & -- & L\\
MYS & G1 & ARDL & 78 & P & 0,149 & 0,909 & -- & --\\
MYS & G1 & NARDL & 78 & P & 1,128 & 0,624 & -- & --\\
SGP & G3 & ARDL & 78 & P & 6,741 & 0,082 & -- & --\\
SGP & G3 & NARDL & 78 & P & 4,384 & 0,138 & -- & --\\
SGP & G5 & ARDL & 78 & P & -0,663 & 0,562 & -- & --\\
SGP & G5 & NARDL & 78 & P & 0,366 & 0,884 & -- & --\\
SGP & G6 & ARDL & 78 & P & -0,766 & 0,097 & -- & L\\
SGP & G6 & NARDL & 78 & P & -0,210 & 0,817 & -- & L\\
SGP & G2 & ARDL & 78 & P & 2,333 & 0,410 & -- & --\\
SGP & G2 & NARDL & 78 & P & -0,273 & 0,961 & -- & --\\
SGP & G7 & ARDL & 78 & P & -0,268 & 0,662 & -- & L\\
SGP & G7 & NARDL & 78 & P & -0,033 & 0,969 & -- & L\\
SGP & G4 & ARDL & 78 & P & 0,219 & 0,794 & -- & --\\
SGP & G4 & NARDL & 78 & P & -2,065 & 0,022 & -- & L\\
SGP & G1 & ARDL & 78 & P & -0,293 & 0,780 & -- & --\\
SGP & G1 & NARDL & 78 & P & 0,500 & 0,792 & -- & --\\
THA & G3 & ARDL & 78 & P & 1,864 & 0,139 & -- & L\\
THA & G3 & NARDL & 78 & P & 5,228 & 0,035 & -- & L\\
THA & G5 & ARDL & 78 & P & -1,181 & 0,303 & -- & --\\
THA & G5 & NARDL & 78 & P & 1,999 & 0,413 & -- & L\\
THA & G6 & ARDL & 78 & P & 1,262 & 0,101 & -- & --\\
THA & G6 & NARDL & 78 & P & 3,030 & 0,050 & -- & --\\
THA & G2 & ARDL & 78 & P & 2,405 & 0,013 & -- & L\\
THA & G2 & NARDL & 78 & P & 0,087 & 0,965 & -- & L\\
THA & G7 & ARDL & 78 & P & -0,811 & 0,142 & -- & --\\
THA & G7 & NARDL & 78 & P & -0,378 & 0,739 & -- & --\\
THA & G4 & ARDL & 78 & P & 0,860 & 0,100 & -- & L+D\\
THA & G4 & NARDL & 78 & P & -0,961 & 0,375 & -- & L\\
THA & G1 & ARDL & 78 & P & 0,883 & 0,284 & -- & --\\
THA & G1 & NARDL & 78 & P & 0,555 & 0,733 & -- & --\\
USA & G3 & ARDL & 75 & P & -0,784 & 0,049 & -- & L+D\\
USA & G3 & NARDL & 75 & P & 0,594 & 0,308 & -- & L\\
USA & G5 & ARDL & 75 & P & 0,392 & 0,910 & -- & --\\
USA & G5 & NARDL & 75 & P & 0,258 & 0,961 & -- & --\\
USA & G6 & ARDL & 75 & P & 1,215 & 0,472 & -- & --\\
USA & G6 & NARDL & 75 & P & 3,305 & \(<0{,}001\) & -- & L+D\\
USA & G2 & ARDL & 75 & P & 0,502 & 0,857 & -- & --\\
USA & G2 & NARDL & 75 & P & 5,114 & \(<0{,}001\) & -- & L+D\\
USA & G7 & ARDL & 75 & P & 10,099 & 0,146 & -- & --\\
USA & G7 & NARDL & 75 & P & 11,962 & 0,027 & -- & --\\
USA & G4 & ARDL & 75 & P & 1,918 & 0,168 & -- & --\\
USA & G4 & NARDL & 75 & P & 3,063 & \(<0{,}001\) & -- & L\\
USA & G1 & ARDL & 75 & P & 0,384 & 0,705 & -- & --\\
USA & G1 & NARDL & 75 & P & 0,643 & 0,699 & -- & --\\
VNM & G3 & ARDL & 74 & P & 6,934 & 0,001 & -- & L+D\\
VNM & G3 & NARDL & 74 & P & 8,952 & 0,002 & -- & L\\
VNM & G5 & ARDL & 74 & P & 3,639 & 0,132 & -- & --\\
VNM & G5 & NARDL & 74 & P & -1,052 & 0,640 & -- & L\\
VNM & G6 & ARDL & 74 & P & 3,707 & 0,048 & -- & L\\
VNM & G6 & NARDL & 74 & P & 4,373 & 0,110 & -- & L\\
VNM & G2 & ARDL & 74 & P & 4,466 & 0,001 & -- & L\\
VNM & G2 & NARDL & 74 & P & 3,771 & 0,032 & -- & L\\
VNM & G7 & ARDL & 74 & P & -4,044 & 0,033 & -- & L\\
VNM & G7 & NARDL & 74 & P & -4,365 & 0,133 & -- & L\\
VNM & G4 & ARDL & 74 & P & -0,155 & 0,879 & -- & L\\
VNM & G4 & NARDL & 74 & P & -1,394 & 0,455 & -- & --\\
VNM & G1 & ARDL & 74 & P & -0,053 & 0,960 & -- & --\\
VNM & G1 & NARDL & 74 & P & -1,964 & 0,070 & -- & --\\
WLD & G3 & ARDL & 61 & P & 0,079 & 0,900 & -- & --\\
WLD & G3 & NARDL & 61 & P & -1,335 & 0,304 & -- & --\\
WLD & G5 & ARDL & 61 & P & 3,900 & 0,147 & -- & --\\
WLD & G5 & NARDL & 61 & P & 2,322 & 0,262 & -- & --\\
WLD & G6 & ARDL & 61 & P & -1,378 & \(<0{,}001\) & -- & L\\
WLD & G6 & NARDL & 61 & P & -0,332 & 0,729 & -- & --\\
WLD & G2 & ARDL & 61 & P & 0,785 & 0,226 & -- & L\\
WLD & G2 & NARDL & 61 & P & 2,901 & 0,003 & -- & L\\
WLD & G7 & ARDL & 61 & P & -2,305 & \(<0{,}001\) & -- & L\\
WLD & G7 & NARDL & 61 & P & -0,525 & 0,348 & -- & L\\
WLD & G4 & ARDL & 61 & P & 0,444 & 0,736 & -- & --\\
WLD & G4 & NARDL & 61 & P & 1,531 & 0,158 & -- & L\\
WLD & G1 & ARDL & 61 & P & 1,362 & 0,085 & -- & --\\
WLD & G1 & NARDL & 61 & P & 3,676 & \(<0{,}001\) & -- & L+D\\
KOR & G0 & ARDL & 78 & FX & -0,093 & 0,930 & -- & --\\
KOR & G0 & NARDL & 78 & FX & 3,846 & 0,156 & -- & --\\
THA & G0 & ARDL & 78 & FX & 1,117 & 0,025 & -- & L\\
THA & G0 & NARDL & 78 & FX & 0,553 & 0,552 & -- & L\\
VNM & G0 & ARDL & 78 & FX & 1,336 & 0,236 & -- & --\\
VNM & G0 & NARDL & 78 & FX & 0,091 & 0,949 & -- & --\\
IND & G0 & ARDL & 77 & WPI & 0,191 & 0,874 & -- & --\\
IND & G0 & NARDL & 77 & WPI & -2,095 & 0,339 & -- & --\\
AUS & G3 & ARDL & 79 & L2 & -2,185 & 0,011 & -- & L\\
AUS & G3 & NARDL & 79 & L2 & -3,904 & \(<0{,}001\) & -- & L+D\\
AUS & G5 & ARDL & 79 & L2 & 2,905 & 0,270 & -- & --\\
AUS & G5 & NARDL & 79 & L2 & 5,587 & 0,136 & -- & --\\
AUS & G6 & ARDL & 79 & L2 & 1,442 & 0,056 & -- & L\\
AUS & G6 & NARDL & 79 & L2 & 0,764 & 0,474 & -- & L\\
AUS & G2 & ARDL & 79 & L2 & 0,617 & 0,367 & -- & --\\
AUS & G2 & NARDL & 79 & L2 & 0,928 & 0,348 & -- & --\\
AUS & G7 & ARDL & 79 & L2 & 3,662 & 0,030 & -- & L\\
AUS & G7 & NARDL & 79 & L2 & 4,368 & 0,068 & -- & L\\
AUS & G4 & ARDL & 79 & L2 & 1,522 & 0,024 & -- & L\\
AUS & G4 & NARDL & 79 & L2 & 1,875 & 0,050 & -- & L\\
AUS & G1 & ARDL & 79 & L2 & 0,484 & 0,710 & -- & --\\
AUS & G1 & NARDL & 79 & L2 & 2,465 & 0,136 & -- & --\\
AUS & G0 & ARDL & 79 & L2 & 1,099 & 0,060 & -- & L\\
AUS & G0 & NARDL & 79 & L2 & 1,818 & 0,038 & -- & L+D\\
CHN & G3 & ARDL & 80 & L2 & -0,569 & 0,173 & -- & --\\
CHN & G3 & NARDL & 80 & L2 & -1,643 & 0,014 & -- & --\\
CHN & G5 & ARDL & 80 & L2 & 1,081 & 0,318 & -- & L\\
CHN & G5 & NARDL & 80 & L2 & 0,540 & 0,752 & -- & L\\
CHN & G6 & ARDL & 80 & L2 & -1,325 & 0,023 & -- & L\\
CHN & G6 & NARDL & 80 & L2 & -3,051 & \(<0{,}001\) & -- & L\\
CHN & G2 & ARDL & 80 & L2 & 0,281 & 0,784 & -- & --\\
CHN & G2 & NARDL & 80 & L2 & -1,750 & 0,180 & -- & --\\
CHN & G7 & ARDL & 80 & L2 & -1,376 & 0,148 & -- & --\\
CHN & G7 & NARDL & 80 & L2 & -1,995 & 0,126 & -- & L\\
CHN & G4 & ARDL & 80 & L2 & -2,384 & 0,035 & -- & --\\
CHN & G4 & NARDL & 80 & L2 & -5,804 & \(<0{,}001\) & -- & --\\
CHN & G1 & ARDL & 80 & L2 & 1,227 & 0,635 & -- & --\\
CHN & G1 & NARDL & 80 & L2 & 26,044 & 0,531 & -- & --\\
CHN & G0 & ARDL & 80 & L2 & -1,614 & 0,090 & -- & --\\
CHN & G0 & NARDL & 80 & L2 & -1,616 & 0,313 & -- & --\\
IND & G3 & ARDL & 51 & L2 & 1,083 & 0,833 & -- & --\\
IND & G3 & NARDL & 51 & L2 & -3,970 & 0,292 & -- & L\\
IND & G5 & ARDL & 51 & L2 & 5,581 & 0,435 & -- & --\\
IND & G5 & NARDL & 51 & L2 & 6,162 & 0,443 & -- & --\\
IND & G6 & ARDL & 51 & L2 & 0,709 & 0,896 & -- & --\\
IND & G6 & NARDL & 51 & L2 & -5,063 & 0,141 & -- & L\\
IND & G2 & ARDL & 51 & L2 & -0,599 & 0,907 & -- & --\\
IND & G2 & NARDL & 51 & L2 & -3,282 & 0,442 & -- & L\\
IND & G7 & ARDL & 51 & L2 & 1,226 & 0,856 & -- & --\\
IND & G7 & NARDL & 51 & L2 & -8,131 & 0,108 & -- & L\\
IND & G4 & ARDL & 51 & L2 & 2,231 & 0,672 & -- & --\\
IND & G4 & NARDL & 51 & L2 & -3,881 & 0,250 & -- & --\\
IND & G1 & ARDL & 51 & L2 & -2,090 & 0,512 & -- & --\\
IND & G1 & NARDL & 51 & L2 & -2,657 & 0,346 & -- & --\\
IND & G0 & ARDL & 51 & L2 & -0,965 & 0,809 & -- & --\\
IND & G0 & NARDL & 51 & L2 & -4,644 & 0,099 & -- & --\\
JPN & G3 & ARDL & 80 & L2 & -0,598 & 0,298 & -- & L\\
JPN & G3 & NARDL & 80 & L2 & -2,018 & \(<0{,}001\) & -- & L\\
JPN & G5 & ARDL & 80 & L2 & -1,975 & 0,305 & -- & --\\
JPN & G5 & NARDL & 80 & L2 & -3,883 & 0,039 & -- & --\\
JPN & G6 & ARDL & 80 & L2 & -1,918 & \(<0{,}001\) & -- & L\\
JPN & G6 & NARDL & 80 & L2 & -1,031 & 0,245 & -- & L+D\\
JPN & G2 & ARDL & 80 & L2 & -0,977 & 0,244 & -- & L+D\\
JPN & G2 & NARDL & 80 & L2 & -1,623 & 0,166 & -- & L\\
JPN & G7 & ARDL & 80 & L2 & -2,267 & 0,002 & -- & L\\
JPN & G7 & NARDL & 80 & L2 & -0,482 & 0,648 & -- & L+D\\
JPN & G4 & ARDL & 80 & L2 & -2,152 & 0,004 & -- & --\\
JPN & G4 & NARDL & 80 & L2 & -2,345 & 0,031 & -- & --\\
JPN & G1 & ARDL & 80 & L2 & -1,490 & 0,164 & -- & --\\
JPN & G1 & NARDL & 80 & L2 & -1,869 & 0,155 & -- & --\\
JPN & G0 & ARDL & 80 & L2 & -2,472 & 0,002 & -- & --\\
JPN & G0 & NARDL & 80 & L2 & -2,832 & 0,007 & -- & --\\
KOR & G3 & ARDL & 80 & L2 & 1,286 & 0,029 & -- & --\\
KOR & G3 & NARDL & 80 & L2 & 2,287 & 0,009 & -- & --\\
KOR & G5 & ARDL & 80 & L2 & 0,989 & 0,474 & -- & --\\
KOR & G5 & NARDL & 80 & L2 & 4,392 & 0,060 & -- & L\\
KOR & G6 & ARDL & 80 & L2 & -0,538 & 0,508 & -- & --\\
KOR & G6 & NARDL & 80 & L2 & 2,158 & 0,171 & -- & --\\
KOR & G2 & ARDL & 80 & L2 & 5,675 & 0,053 & -- & --\\
KOR & G2 & NARDL & 80 & L2 & 11,736 & 0,057 & -- & --\\
KOR & G7 & ARDL & 80 & L2 & -0,716 & 0,469 & -- & --\\
KOR & G7 & NARDL & 80 & L2 & 2,647 & 0,176 & -- & --\\
KOR & G4 & ARDL & 80 & L2 & 1,992 & 0,460 & -- & --\\
KOR & G4 & NARDL & 80 & L2 & 10,901 & 0,079 & -- & --\\
KOR & G1 & ARDL & 80 & L2 & 1,714 & 0,031 & -- & --\\
KOR & G1 & NARDL & 80 & L2 & 3,624 & 0,007 & -- & L\\
KOR & G0 & ARDL & 80 & L2 & -0,043 & 0,944 & -- & --\\
KOR & G0 & NARDL & 80 & L2 & 1,496 & 0,177 & -- & L\\
MYS & G3 & ARDL & 80 & L2 & -0,051 & 0,964 & -- & --\\
MYS & G3 & NARDL & 80 & L2 & 3,665 & 0,109 & -- & --\\
MYS & G5 & ARDL & 80 & L2 & 0,207 & 0,705 & -- & L+D\\
MYS & G5 & NARDL & 80 & L2 & 0,530 & 0,616 & -- & L+D\\
MYS & G6 & ARDL & 80 & L2 & -1,284 & 0,056 & -- & L+D\\
MYS & G6 & NARDL & 80 & L2 & 1,201 & 0,273 & -- & L+D\\
MYS & G2 & ARDL & 80 & L2 & -1,533 & 0,475 & -- & --\\
MYS & G2 & NARDL & 80 & L2 & 4,639 & 0,044 & -- & --\\
MYS & G7 & ARDL & 80 & L2 & -1,272 & 0,132 & -- & L+D\\
MYS & G7 & NARDL & 80 & L2 & 0,724 & 0,632 & -- & L+D\\
MYS & G4 & ARDL & 80 & L2 & 0,412 & 0,572 & -- & --\\
MYS & G4 & NARDL & 80 & L2 & 1,285 & 0,396 & -- & --\\
MYS & G1 & ARDL & 80 & L2 & 0,144 & 0,883 & -- & --\\
MYS & G1 & NARDL & 80 & L2 & 1,785 & 0,266 & -- & --\\
MYS & G0 & ARDL & 80 & L2 & -0,713 & 0,085 & -- & L\\
MYS & G0 & NARDL & 80 & L2 & 1,356 & 0,021 & -- & L+D\\
SGP & G3 & ARDL & 80 & L2 & 3,033 & 0,106 & -- & --\\
SGP & G3 & NARDL & 80 & L2 & 5,321 & 0,039 & -- & --\\
SGP & G5 & ARDL & 80 & L2 & -0,446 & 0,680 & -- & --\\
SGP & G5 & NARDL & 80 & L2 & -1,138 & 0,518 & -- & --\\
SGP & G6 & ARDL & 80 & L2 & -1,185 & 0,031 & -- & --\\
SGP & G6 & NARDL & 80 & L2 & -0,429 & 0,625 & -- & --\\
SGP & G2 & ARDL & 80 & L2 & 2,773 & 0,321 & -- & --\\
SGP & G2 & NARDL & 80 & L2 & -2,652 & 0,457 & -- & L\\
SGP & G7 & ARDL & 80 & L2 & -1,380 & 0,019 & -- & --\\
SGP & G7 & NARDL & 80 & L2 & -0,523 & 0,576 & -- & --\\
SGP & G4 & ARDL & 80 & L2 & 0,400 & 0,651 & -- & --\\
SGP & G4 & NARDL & 80 & L2 & -2,152 & 0,017 & -- & L\\
SGP & G1 & ARDL & 80 & L2 & -0,252 & 0,857 & -- & --\\
SGP & G1 & NARDL & 80 & L2 & -1,741 & 0,308 & -- & --\\
SGP & G0 & ARDL & 80 & L2 & -0,640 & 0,191 & -- & --\\
SGP & G0 & NARDL & 80 & L2 & -1,522 & 0,038 & -- & --\\
THA & G3 & ARDL & 80 & L2 & 1,774 & 0,175 & -- & L+D\\
THA & G3 & NARDL & 80 & L2 & 3,248 & 0,192 & -- & L+D\\
THA & G5 & ARDL & 80 & L2 & -1,960 & 0,099 & -- & --\\
THA & G5 & NARDL & 80 & L2 & -0,517 & 0,817 & -- & --\\
THA & G6 & ARDL & 80 & L2 & 1,008 & 0,208 & -- & --\\
THA & G6 & NARDL & 80 & L2 & 0,377 & 0,802 & -- & --\\
THA & G2 & ARDL & 80 & L2 & 2,466 & 0,026 & -- & L\\
THA & G2 & NARDL & 80 & L2 & -0,963 & 0,715 & -- & --\\
THA & G7 & ARDL & 80 & L2 & -0,561 & 0,414 & -- & --\\
THA & G7 & NARDL & 80 & L2 & -0,140 & 0,919 & -- & --\\
THA & G4 & ARDL & 80 & L2 & 0,969 & 0,122 & -- & L\\
THA & G4 & NARDL & 80 & L2 & -0,197 & 0,860 & -- & L\\
THA & G1 & ARDL & 80 & L2 & 0,191 & 0,829 & -- & --\\
THA & G1 & NARDL & 80 & L2 & -1,896 & 0,217 & -- & --\\
THA & G0 & ARDL & 80 & L2 & 1,144 & 0,044 & -- & L\\
THA & G0 & NARDL & 80 & L2 & -0,400 & 0,684 & -- & L\\
USA & G3 & ARDL & 77 & L2 & -0,823 & 0,016 & -- & L+D\\
USA & G3 & NARDL & 77 & L2 & 0,573 & 0,288 & -- & L+D\\
USA & G5 & ARDL & 77 & L2 & 4,163 & 0,257 & -- & --\\
USA & G5 & NARDL & 77 & L2 & 2,780 & 0,597 & -- & --\\
USA & G6 & ARDL & 77 & L2 & -0,353 & 0,554 & -- & L\\
USA & G6 & NARDL & 77 & L2 & 1,659 & 0,027 & -- & L\\
USA & G2 & ARDL & 77 & L2 & 0,871 & 0,526 & -- & --\\
USA & G2 & NARDL & 77 & L2 & 5,043 & \(<0{,}001\) & -- & L+D\\
USA & G7 & ARDL & 77 & L2 & 1,924 & 0,611 & -- & --\\
USA & G7 & NARDL & 77 & L2 & 1,977 & 0,746 & -- & --\\
USA & G4 & ARDL & 77 & L2 & 0,080 & 0,914 & -- & --\\
USA & G4 & NARDL & 77 & L2 & 1,668 & 0,031 & -- & L\\
USA & G1 & ARDL & 77 & L2 & 1,644 & 0,186 & -- & --\\
USA & G1 & NARDL & 77 & L2 & 3,403 & 0,088 & -- & --\\
USA & G0 & ARDL & 77 & L2 & 1,080 & 0,187 & -- & --\\
USA & G0 & NARDL & 77 & L2 & 3,284 & 0,005 & -- & L\\
VNM & G3 & ARDL & 76 & L2 & 5,938 & 0,006 & -- & L+D\\
VNM & G3 & NARDL & 76 & L2 & 7,636 & 0,013 & -- & --\\
VNM & G5 & ARDL & 76 & L2 & 2,771 & 0,125 & -- & --\\
VNM & G5 & NARDL & 76 & L2 & -0,477 & 0,805 & -- & L\\
VNM & G6 & ARDL & 76 & L2 & 2,254 & 0,232 & -- & L\\
VNM & G6 & NARDL & 76 & L2 & 2,601 & 0,372 & -- & --\\
VNM & G2 & ARDL & 76 & L2 & 4,033 & 0,006 & -- & L\\
VNM & G2 & NARDL & 76 & L2 & 4,771 & 0,028 & -- & L\\
VNM & G7 & ARDL & 76 & L2 & -3,782 & 0,039 & -- & L\\
VNM & G7 & NARDL & 76 & L2 & -3,585 & 0,199 & -- & L\\
VNM & G4 & ARDL & 76 & L2 & -0,856 & 0,418 & -- & L\\
VNM & G4 & NARDL & 76 & L2 & -2,427 & 0,105 & -- & L\\
VNM & G1 & ARDL & 76 & L2 & 1,312 & 0,361 & -- & --\\
VNM & G1 & NARDL & 76 & L2 & -0,865 & 0,534 & -- & --\\
VNM & G0 & ARDL & 76 & L2 & 0,739 & 0,365 & -- & L\\
VNM & G0 & NARDL & 76 & L2 & 0,236 & 0,850 & -- & L\\
WLD & G3 & ARDL & 63 & L2 & -0,279 & 0,657 & -- & --\\
WLD & G3 & NARDL & 63 & L2 & 0,292 & 0,767 & -- & --\\
WLD & G5 & ARDL & 63 & L2 & 0,877 & 0,511 & -- & --\\
WLD & G5 & NARDL & 63 & L2 & 1,422 & 0,497 & -- & --\\
WLD & G6 & ARDL & 63 & L2 & -1,536 & \(<0{,}001\) & -- & L+D\\
WLD & G6 & NARDL & 63 & L2 & 0,681 & 0,281 & -- & L+D\\
WLD & G2 & ARDL & 63 & L2 & 1,247 & 0,319 & -- & --\\
WLD & G2 & NARDL & 63 & L2 & 3,494 & 0,051 & -- & --\\
WLD & G7 & ARDL & 63 & L2 & -2,222 & \(<0{,}001\) & -- & L\\
WLD & G7 & NARDL & 63 & L2 & -0,715 & 0,288 & -- & L+D\\
WLD & G4 & ARDL & 63 & L2 & -1,886 & 0,037 & -- & --\\
WLD & G4 & NARDL & 63 & L2 & 1,008 & 0,463 & -- & --\\
WLD & G1 & ARDL & 63 & L2 & 0,454 & 0,537 & -- & --\\
WLD & G1 & NARDL & 63 & L2 & 2,675 & 0,002 & -- & L+D\\
WLD & G0 & ARDL & 63 & L2 & -0,728 & 0,128 & -- & --\\
WLD & G0 & NARDL & 63 & L2 & 1,845 & 0,044 & -- & --\\
WLD & G0 & ARDL & 77 & W06 & 1,717 & 0,422 & -- & --\\
WLD & G0 & NARDL & 77 & W06 & 6,321 & 0,293 & -- & --\\
\end{longtable}
}

\clearpage\section{Glosarium interpretasi}
\begin{longtable}{p{.24\linewidth}p{.65\linewidth}}
\toprule Istilah & Arti dalam penelitian\\\midrule\endfirsthead
\toprule Istilah & Arti dalam penelitian\\\midrule\endhead
Outcome $y$ & Log impor Indonesia dibagi ekspor pada cakupan tertentu; turun berarti perbaikan relatif.\\
Kurs $q$ & Log indeks kurs riil dengan orientasi naik sebagai apresiasi.\\
Partial sums & Akumulasi perubahan positif dan negatif; bukan ukuran volatilitas.\\
Koefisien $\beta_-$ & Hubungan jangka panjang dengan partial sum depresiasi; tanda harus dikalikan perubahan kurs negatif.\\
Error correction $\rho$ & Parameter penyesuaian terhadap deviasi level. Negatif saja belum membuktikan hubungan level sah.\\
Bounds F dan t & Pemeriksaan hubungan level memakai distribusi kritis khusus, bukan sekadar p-value regresi biasa.\\
Joint-X & Uji tambahan komponen level regressor untuk menilai bentuk hubungan yang didukung.\\
CI pointwise & Interval pada setiap horizon, bukan jaminan simultan seluruh kurva.\\
Nominal pascaseleksi & Probabilitas dan interval belum mengoreksi proses pemilihan lag dengan AIC.\\
BH & Koreksi pengujian berganda dalam keluarga yang ditetapkan; penerapan di sini eksploratif.\\
Pola J titik & Klasifikasi operasional pemburukan awal dan perbaikan akhir pada estimasi titik.\\
Horizon 24 & Enam tahun sejak guncangan dalam simulasi dinamis, tidak otomatis long-run limit.\\
Tidak signifikan & Bukti belum cukup menolak efek nol di bawah asumsi, bukan bukti efek benar-benar nol.\\
Mekanisme & Penjelasan seperti harga atau volume yang memerlukan pengukuran dan identifikasi lebih lanjut.\\
Integritas pengolahan & Konsistensi asal, aritmetika, kunci waktu dan numerik; berbeda dari validitas kausal atau keseragaman konsep sumber.\\
\bottomrule\end{longtable}

\section{Jejak reproduksi}
Data perdagangan mentah berada pada \path{input/raw/trademap_hs6_20261117/}; referensi klasifikasi dan paper pada \path{input/reference/}. Panel antara dibentuk oleh skrip persiapan dan tidak diperlakukan sebagai data mentah. Master Stata menjalankan model serta diagnostik; hasil pelaporan diperiksa melalui ekuivalensi ECM--level dan prediksi OLS independen.

Pembuatan paper menggunakan \path{script/12_build_research_paper.py}, sumber \path{paper.tex}, dan folder \path{sections/}. Generator menarik angka dan grafik dari hasil yang telah selesai, tanpa mengubah data atau mengestimasi ulang model. \path{paper_validation.json} memuat checksum sumber, jumlah model yang dilaporkan, dan hasil pemeriksaan PDF. Source LaTeX disertakan agar uraian dapat diperbarui ketika penelitian berkembang.
